\section{All-ideal reciprocity and the surviving block obstruction}
\label{wc:details}
The proof uses no new zero-free input. Its external analytic ingredient
is the pinned BGL large sieve and the powered estimates established in
its proof; elementary ideal counting and the already restored Mellin
identity supply the remaining interfaces. All constants may depend on
fixed $K$, $S$, $\nu$ and declared small exponents, never on $U$.

\subsection{Exact recombination retains the primitive corrections}
On $\operatorname{Re}s>1$, start from the full defining series, with
its original zero masks:
\[
 F_{u,d}(s)=\sum_{c\ \mathrm{squarefree}\atop(c,T(u,d))=1}
                 \nu(c)\chi_c(u)(Nc)^{-s}.
\]
Its primitive ratio involves $L(s,\psi_u)$, the independently primitive
$L(2s,\psi_{u,2})$ and every deleted Euler factor. Expanding this
\emph{restored} series gives the exact equality
\begin{align}\label{wc:recombine}
 &\sum_{d\ \mathrm{squarefree}\atop(d,S)=1}
 \beta_U(d)\nu(d)\chi_d(u)(Nd)^{-s}F_{u,d}(s)\notag\\
 &\qquad=\sum_{m\ \mathrm{squarefree}\atop(m,S)=1}
               A_U(m)\nu(m)\chi_m(u)(Nm)^{-s}.
\end{align}
Indeed $(c,d)=1$ makes $m=cd$ squarefree, with $d\mid m$. If $m$
meets $u$, both original symbols retain their zero. Only on the
absolutely convergent line is this infinite series used. Mellin inversion
there gives the finite smoothed polynomial~\eqref{wc:polynomial}.
Nothing evaluates~\eqref{wc:recombine} as a Dirichlet series on
$\operatorname{Re}s=1/2+\delta$ or interchanges a cutoff limit with
continuation across unknown zeros.

For squarefree $m$, a factorization $ab=d\mid m$ assigns each prime
to $a$, to $b$, or to neither. The actual inequalities $Na,Nb\le U$
select a subset of these assignments. Thus
\[
 |A_U(m)|\le3^{\omega(m)}\ll_\eta(Nm)^\eta.
\]
The last divisor bound follows by choosing a fixed threshold so
$3\le(Np)^\eta$ beyond it, then absorbing the finitely many smaller
prime ideals in a constant. When $Nm\le U$ all assignments occur,
so $A_U(m)=\prod_{p\mid m}(1-1-1)=\mu(m)$; no such simplification
is made for the length-$D$ columns in general. With
$\lambda_m=D^{-1/2}A_U(m)\nu(m)W(Nm/D)$, ideal counting and a
smaller internal $\eta$ give
\begin{equation}\label{wc:l2}
 \sum_m|\lambda_m|^2\ll_\eta\|W\|_\infty^2D^\eta,
 \qquad \operatorname{supp}\lambda\subseteq\{D\le Nm\le2D\}.
\end{equation}

\subsection{The powered squarefree input and all-ideal pivot}
Let $\vartheta_a$ be BGL's canonical primitive sixth-order family.
Their Section~3 defines $B_1^j$ using squarefree row/column ideals and
kernel $\vartheta_a(m)^j$, while $B_2^j$ uses \emph{all} ideal rows
and kernel $\vartheta_m(a)^j$. Theorem~1.3 states the squarefree bound
for $j=1$; its Section~3.2 proves the estimates simultaneously for
$1\le j\le5$, because the recursion in Lemma~3.4 uses $j\ell$.
In particular the two inputs needed here are
\begin{equation}\label{wc:powered-input}
 B_1^j(M,D)\ll_\eta(MD)^\eta
                 \{M+D+(MD)^{2/3}\},\qquad j=1,2.
\end{equation}
To see the final step in the published recursion explicitly, its
$E_{4/3}$ estimate is $M^{4/3}+D+(MD)^{2/3}$.
If $M\le D$, its first term is at most $(MD)^{2/3}$; if $M>D$,
use reciprocity/duality (3.2) and the estimate with $M,D$ reversed.
The finite set of powers is treated together; no inference from a
fixed conductor count to an arbitrary primitive-character sieve occurs.

Write each ideal row uniquely as
\[
 a=a_1a_2^2a_3^3a_4^4a_5^5a_6^6,
\]
where $a_1,\ldots,a_5$ are pairwise-coprime squarefree ideals and
$a_6$ is arbitrary. For valuation $6q+r$, $0\le r\le5$, place the
prime in $a_r$ if $r>0$ and its power $q$ in $a_6$.
Dyadically let $A_j\le Na_j<2A_j$, with $A_j\ge1$, and put
$R=\prod_{j=1}^6A_j$. A nonempty block satisfies
$\prod_jA_j^j\le X$ when $Na\le X$.
Choose $\ell=1$ if $A_1\ge A_2$, and $\ell=2$ otherwise; put
$M=A_\ell$. Keep every other ideal fixed. Evaluation of the single
\emph{column} character gives
\[
 \vartheta_m(a)=\vartheta_m(a_\ell)^\ell
                  \vartheta_m(a/a_\ell^\ell).
\]
The second factor is a bounded column multiplier, including its original
zero if $m$ meets any fixed factor. It may depend on all the fixed
ideals but not on the pivot row. Pairwise-coprimality and the exact row
norm restriction merely select a subset of the positive pivot-row sum.
Apply~\eqref{wc:powered-input} to the first factor. There are
$O_K(R/M)$ choices of the other ideals by ideal counting. The block
energy is bounded by
\begin{equation}\label{wc:block}
 (XD)^\eta\{R+DR/M+D^{2/3}R/M^{1/3}\}\sum_m|\lambda_m|^2.
\end{equation}

All three monomial bounds used to sum these blocks are elementary.
First $R\le\prod_jA_j^j\le X$.
If $M=A_1$, then
\[
 (R/M)^3=A_2^3\cdots A_6^3
 \le A_1A_2^2A_3^3A_4^4A_5^5A_6^6\le X,
\]
since $A_2\le A_1$ and $j\ge3$ for the remaining factors.
If $M=A_2$, use $A_1\le A_2$ in the same comparison.
Finally,
\[
 \frac{R^3}{M(\prod_jA_j^j)^2}
 =\frac{A_1}{MA_2A_3^3A_4^5A_5^7A_6^9}\le1.
\]
Thus $R/M\le X^{1/3}$ and $R/M^{1/3}\le X^{2/3}$.
Insert these in~\eqref{wc:block} and sum the logarithmically many
blocks. This proves the all-ideal estimate used in Theorem~\ref{wc:uniform}.
It is precisely the mechanism of BGL Lemma~3.4 and (3.6)--(3.9),
with the final powered bound inserted. The simpler fixed-squarefull-row
argument would leave the weaker $D\sqrt X$ term.

\subsection{Literal rows, reciprocal orientation and principal count}
Let $a=(u)_S$ denote the outside-$S$ ideal. The finite $S$-unit
comparison gives $\kappa_u=\eta_u\vartheta_a$ as primitive
characters, with $\eta_u$ in a finite fixed set.
For $(a,m)=1$, BGL reciprocity therefore gives
\begin{equation}\label{wc:literal-reciprocity}
 \chi_m(u)=\eta_u(m)C([a],[m])\vartheta_m(a).
\end{equation}
The notation is deliberate: the left side is the project's original
$(u/m)_6$, not BGL's canonical character indexed by $m$.
When $(a,m)>1$, the left side is zero. The right side is also zero,
because squarefree $m$ has outside-$S$ primitive conductor exactly $m$.
Equation~\eqref{wc:literal-reciprocity} hence extends to every pair.
Using $\vartheta_a(m)$ on noncoprime pairs instead would be false:
primitive $\vartheta_{b^6}$ has restored primes of $b$, whereas the
original symbol must still vanish there.

Partition rows by $\eta_u$ and $[a]$, both finite. On each part the
factor $\eta_u(m)C([a],[m])$ is a fixed modulus-one column multiplier.
A positive row sum over that part is bounded by the full $B_2^1$ row
sum. For each $S$-supported ideal $b$ and outside-$S$ ideal $a$, at
most six actual elements generate $ab$. Thus the required literal energy
is bounded by the all-ideal norm summed over $b$, with $X=H/Nb$.
The geometric products
\[
 \sum_{b\ S\text{-supported}}(Nb)^{-q}
   =\prod_{p\in S}(1-(Np)^{-q})^{-1}<\infty,
 \qquad q=1,\ 1/3,\ 2/3,
\]
transfer the three terms without losing any unit or $S$-power row.
Together with~\eqref{wc:l2}, this proves the claimed $I$ energy.

If $\psi_u$ is principal, $\nu$ is unramified outside $S$, and tame
Kummer inertia at $v_p(u)$ has order $6/\gcd(6,v_p(u))$.
It must be trivial, so each such valuation is divisible by $6$.
Writing $(u)=ba_6^6$ with $b$ supported in $S$, ideal counting gives
\[
 \#\{0<Nu\le H:\psi_u=1\}
 \ll_K H^{1/6}\prod_{p\in S}(1-(Np)^{-1/6})^{-1}.
\]
This is only a necessary condition, sufficient for an upper count.
The exact residue formula~\eqref{wc:principal} and
$|\widehat W(1)|\le\|W\|_\infty$ then give the principal energy
bound used above. No sign or positivity of that residue is assumed.

\subsection{What the envelope still fails to prove}
The $DH^{1/3}$ term is still larger than $H(HD)^\varepsilon$ at
$H=D^{3/4}$. The monomial comparison itself is saturated by
$A_1=A_2=X^{1/3}$ and $A_3=\cdots=A_6=1$.
This proves sharpness of that \emph{particular comparison}, not an
energy lower bound for actual arithmetic coefficients. In particular,
no conclusion that the signed $A_U$ cannot cancel is justified.
The next load-bearing input must exploit those actual truncated M\"obius
coefficients inside the mixed moment, or establish the required signed
type-II minus type-I covariance. The smaller-scale short term, principal
contribution, full arithmetic determinant identification and GRH remain
separate obligations.

Exact finite checks use the split Eisenstein primes generated by
$3+\omega$, $4+\omega$, $5+2\omega$, $6+\omega$, of norms
$7,13,19,31$. Their actual residue symbols are computed in
$\mathbb Q(\zeta_6)$, retaining sixth-power zeros. Coefficient
recombination and dyadic inequalities pass; they check finite algebra
and bookkeeping, not the external analytic sieve or the contour theorem.
No new Lean declarations accompany this argument.
